\scanpage{99}
% archive foliation: 98
que j'écris sous forme contragrédiente
\[
\check\kappa_i^{\otimes2}(\check\delta^{\otimes2})
=\check\nu(\check\delta^{\otimes2})\wedge\omega_i,
\]
en termes des $\psi^\omega(\delta)$, comme
\[
(73)\qquad\psi^\omega(\delta)^{\otimes2}
=\check\nu(\check\delta^{\otimes2})\wedge\omega.
\]
% margin position: right
\margin{\uline{NB} Moyennant (71)\inserted{(et 72)}, les deux
relations (73), pour les deux él. $\omega\in\underline\omega$, sont équivalentes,
\unsure{il suffit} \unsure{sa validité} pour un $\delta$
\deleted{\unsure{l'indiquer}}\illegible{4}.}
On trouve ainsi le

\uline{Lemme} Les épinglages de Legendre \unsure{de la}
donnée $(J,F_t\subset V_J(k),\underline t,\check\nu:\check{\underline t}^{\otimes2}\simeq F_t^{\otimes2}(\underline\omega))$,
s'identifient aux systèmes $(\alpha,(\psi^\omega(\delta))_{\omega\in\underline\omega})$, où
$\alpha:\mu_4^*\simeq\underline\omega$, et où
$(\psi^\omega(\delta))_{\substack{\omega\in\underline\omega\\\delta\in\underline t^*}}\in(F_t)^{\underline\omega\times\underline t^*}$
\deleted{satisfaisant}
une famille indexée par l'ens. \inserted{produit $\underline\omega\times\underline t^*$}
\deleted{\unsure{à deux élé}}\unsure{ments} $\underline\omega=\underline\omega_J$, satisfaisant les conditions
(71) (72) (73).

\deleted{\unsure{Nous}} \unsure{Nous|Pour} maintenant expliciter la
donnée d'une famille à deux indices
$(\psi_i^\omega(\delta))_{\substack{\omega\in\underline\omega\\\delta\in\underline t^*}}$, à val. dans $F_t$, en termes du
plongement
\[
F_t\subset V_J=\operatorname{Ker}(k^J\xrightarrow[\mathrm{somme}]{}k),
\]
on voit donc qu'elle s'exprime comme
une famille à trois indices d'éléments de $k$:
\[
(74)\qquad(\psi_i^\omega(\delta))_{\substack{i\in J\\\omega\in\underline\omega\\\delta\in\underline t^*}},\qquad\psi_i^\omega(\delta)\in k
\]
Le fait que la famille d'éléments \unsure{$\in$} de $k^J$
